% Generated by publication/build_sections.py; do not edit by hand.
\paragraph{Informal verdict.}
\textbf{A --- Complete solution}, with high confidence, for the meaningful range \(n\ge3\).

\paragraph{Lean coverage.}
The Lean development proves the control and steering lemmas and the complete minimax outcome theorem for every n at least 3. The lower bound is the necessary domain repair because the printed last-mover payoff is undefined for n=1,2.

\subsection{Audited informal answer}
\subsubsection{1. Faithful restatement and a domain issue}

A position is a finite collection of sticks of positive integer lengths. A legal move replaces one stick of length \(m\) by two sticks of distinct positive integer lengths \(a<b\) satisfying \(a+b=m\). Hence a stick is splittable exactly when \(m\ge 3\).

Player \(A\) makes the first move, then \(B\), and they alternate until every remaining stick has length \(1\) or \(2\). If the terminal position contains \(x\) sticks of length \(1\) and \(y\) sticks of length \(2\), then:

\begin{itemize}
\item if \(x>y\), the last mover wins;
\item if \(x<y\), the last mover loses;
\item if \(x=y\), the game is drawn.
\end{itemize}

The question asks whether, under optimal play, the outcome depends only on \(n\bmod 3\) in the stated way. This is the complete statement of Q728 in the canonical PDF. 

There is one genuine omission: the PDF does not explicitly assume \(n\ge 3\). For \(n=1,2\), \(A\) has no legal first move and there is no ``last player'', so the payoff rule is undefined. I therefore solve the intended problem for \(n\ge 3\).

\textbf{Theorem.} For every integer \(n\ge 3\):
\[
\begin{cases}
A\text{ wins},&n\equiv2\pmod3,\\[2mm]
B\text{ wins},&n\equiv1\pmod3,\\[2mm]
\text{the game is drawn},&n\equiv0\pmod3.
\end{cases}
\]
Thus the proposed assertion is true after inserting the necessary hypothesis \(n\ge3\).

The proof is constructive and gives explicit optimal strategies.

\medskip\hrule\medskip

\subsubsection{2. Elementary terminal identities}

Let a terminal position contain \(x\) pieces of length \(1\) and \(y\) pieces of length \(2\). Conservation of total length gives

\[
x+2y=n.
\]

Consequently,

\[
x+y=n-y.
\]

After \(A\)'s opening move there are exactly two pieces, and \(B\) is to move. Since every subsequent move increases the number of pieces by one, the number of moves remaining after the opening is

\[
R=(x+y)-2=n-y-2.
\]

Therefore:

\begin{itemize}
\item \(B\) is the overall last mover exactly when \(R\) is odd;
\item \(A\) is the overall last mover exactly when \(R\) is even, including \(R=0\), since \(A\) then made the opening and final move.
\end{itemize}

The main task is thus to control the possible terminal value of \(y\).

\medskip\hrule\medskip

\subsection{3. The control interval of a position}

For a positive integer \(m\), define

\[
\beta(m)=
\begin{cases}
0,&m=1,\\
1,&m=2,\\
\lfloor m/3\rfloor,&m\ge3,
\end{cases}
\]

and

\[
\eta(m)=
\begin{cases}
1,&m\ge5\text{ and }m\equiv2\pmod3,\\
0,&\text{otherwise}.
\end{cases}
\]

We call a piece counted by \(\eta\) a \textbf{hot piece}.

For a position \(P=(m_1,\ldots,m_k)\), put

\[
b(P)=\sum_{i=1}^k\beta(m_i),
\qquad
h(P)=\sum_{i=1}^k\eta(m_i),
\]

and define its \textbf{control interval}

\[
J(P)=
\left\{
b(P)+\left\lfloor\frac{h(P)}2\right\rfloor,
b(P)+\left\lceil\frac{h(P)}2\right\rceil
\right\}.
\]

Thus \(J(P)\) is either a singleton or a pair of consecutive integers.

If \(P\) is terminal, every piece is \(1\) or \(2\). Hence \(h(P)=0\), while \(b(P)\) is precisely the number \(y\) of pieces of length \(2\). Therefore

\[
\boxed{P\text{ terminal}\quad\Longrightarrow\quad J(P)={y}.}
\]

This is why \(J(P)\) is useful.

\medskip\hrule\medskip

\subsubsection{4. Effect of one move}

Suppose a piece \(m\) is split into \(a+b=m\), with \(1\le a<b\). Write

\[
\Delta b=\beta(a)+\beta(b)-\beta(m),
\qquad
\Delta h=\eta(a)+\eta(b)-\eta(m).
\]

The following table exhausts all possibilities. In the table, residue classes of the children are unordered.

\begin{center}
\small
\begin{tabular}{llll}
\hline
\(m\bmod3\) & Child residues & Additional condition & \((\Delta b,\Delta h)\) \\
\hline
\(0\) & \(0+0\) & --- & ((0,0)) \\
\(0\) & \(1+2\) & the \(2\bmod3\) child is \(2\) & ((0,0)) \\
\(0\) & \(1+2\) & the \(2\bmod3\) child is at least \(5\) & ((-1,1)) \\
\(1\) & \(0+1\) & --- & ((0,0)) \\
\(1\) & \(2+2\) & exactly one child is \(2\) & ((0,1)) \\
\(1\) & \(2+2\) & both children are at least \(5\) & ((-1,2)) \\
\(2\) & \(1+1\) & --- & ((0,-1)) \\
\(2\) & \(0+2\) & the \(2\bmod3\) child is \(2\) & ((1,-1)) \\
\(2\) & \(0+2\) & the \(2\bmod3\) child is at least \(5\) & ((0,0)) \\
\hline
\end{tabular}
\end{center}

Here a splittable piece congruent to \(2\bmod3\) is necessarily at least \(5\), so it is itself hot.

For completeness, the table follows immediately from

\[
\begin{aligned}
(\beta(3r),\eta(3r))&=(r,0),\\
(\beta(3r+1),\eta(3r+1))&=(r,0),\\
(\beta(2),\eta(2))&=(1,0),\\
(\beta(3r+2),\eta(3r+2))&=(r,1),\qquad r\ge1,
\end{aligned}
\]

together with the possible residue decompositions

\[
\begin{aligned}
0&=0+0=1+2,\\
1&=0+1=2+2,\\
2&=0+2=1+1
\end{aligned}
\qquad\pmod3.
\]

Define the half-integral centre

\[
c(P)=b(P)+\frac{h(P)}2.
\]

The table shows that under every move \(P\to P'\),

\[
c(P')-c(P)\in\left\{-\frac12,0,\frac12\right\}.
\]

Since

\[
J(P)={\lfloor c(P)\rfloor,\lceil c(P)\rceil},
\]

we obtain the following overlap property.

\paragraph{Lemma 1 --- Overlap under every move}

For every legal move \(P\to P'\),

\[
\boxed{J(P)\cap J(P')\ne\varnothing.}
\]

Indeed, two half-integers differing by at most (1/2) have integer hulls sharing at least one integer.

\medskip\hrule\medskip

\subsection{5. The steering lemma}

Fix two consecutive integers

\[
T_t={t,t+1}.
\]

\paragraph{Lemma 2 --- One-move steering}

Suppose \(P\) is nonterminal and

\[
J(P)\cap T_t\ne\varnothing.
\]

Then there is a legal move \(P\to P'\) such that

\[
\boxed{J(P')\subseteq T_t.}
\]

\paragraph{Proof}

Write \(b=b(P)\) and \(h=h(P)\).

\subparagraph{Case 1: \(h\) is odd}

Write \(h=2r+1\). Then

\[
J(P)={b+r,b+r+1}.
\]

Because \(h>0\), there is a hot piece \(m\equiv2\pmod3\), \(m\ge5\). Both of the following splits are legal:

\[
m=1+(m-1),
\qquad
m=2+(m-2).
\]

From the move table:

\begin{itemize}
\item the first split has \((\Delta b,\Delta h)=(0,-1)\), and therefore
\end{itemize}
\[
  J(P')={b+r};
\]
\begin{itemize}
\item the second has \((\Delta b,\Delta h)=(1,-1)\), and therefore
\end{itemize}
\[
  J(P')={b+r+1}.
\]

At least one of the two endpoints of \(J(P)\) belongs to \(T_t\). Choose the corresponding split. The new interval is then a singleton contained in \(T_t\).

\subparagraph{Case 2: \(h>0\) is even}

Write \(h=2r\). Then

\[
J(P)={s},\qquad s=b+r.
\]

Since \(J(P)\cap T_t\ne\varnothing\), either \(s=t\) or \(s=t+1\).

Again choose a hot piece \(m\).

\begin{itemize}
\item Splitting \(m=1+(m-1)\) gives
\end{itemize}
\[
  J(P')={s-1,s}.
\]
\begin{itemize}
\item Splitting \(m=2+(m-2)\) gives
\end{itemize}
\[
  J(P')={s,s+1}.
\]

Thus:

\begin{itemize}
\item if \(s=t\), use (2+\(m-2\));
\item if \(s=t+1\), use (1+\(m-1\)).
\end{itemize}

In either case \(J(P')=T_t\).

\subparagraph{Case 3: \(h=0\)}

Then \(J(P)={b}\), and \(b\in T_t\).

Because \(P\) is nonterminal, it contains a splittable piece \(m\ge3\). Since \(h=0\), no splittable piece can be congruent to \(2\bmod3\). Hence \(m\equiv0\) or \(1\bmod3\).

\begin{itemize}
\item If \(m=3\), split it as \(1+2\).
\item If \(m\equiv0\pmod3\) and \(m\ge6\), split it as
\end{itemize}
\[
  m=2+(m-2).
\]
\begin{itemize}
\item If \(m\equiv1\pmod3\), split it as
\end{itemize}
\[
  m=1+(m-1).
\]

Each of these moves has \((\Delta b,\Delta h)=(0,0)\), so

\[
J(P')=J(P)={b}\subseteq T_t.
\]

This proves the lemma. \(\square\)

\medskip\hrule\medskip

\subsection{6. The interval-control theorem}

The previous two lemmas combine into a robust strategy.

\paragraph{Theorem 3 --- A designated player can control the terminal value}

Let \(P_0\) be a position satisfying

\[
J(P_0)\subseteq T_t={t,t+1}.
\]

Designate either of the two players as the \textbf{controller}. It does not matter whether the controller or the opponent is next to move.

Then the controller has a strategy forcing the terminal number \(y\) of length-\(2\) pieces to satisfy

\[
\boxed{y\in{t,t+1}.}
\]

\paragraph{Proof}

The controller maintains the following invariant:

\begin{itemize}
\item immediately after every controller move,
\end{itemize}
\[
  J(P)\subseteq T_t;
\]
\begin{itemize}
\item immediately after every opponent move,
\end{itemize}
\[
  J(P)\cap T_t\ne\varnothing.
\]

Initially \(J(P_0)\subseteq T_t\).

Suppose the controller has just moved and \(J(P)\subseteq T_t\). If the opponent makes a move \(P\to Q\), Lemma 1 gives

\[
J(P)\cap J(Q)\ne\varnothing.
\]

Since \(J(P)\subseteq T_t\), it follows that

\[
J(Q)\cap T_t\ne\varnothing.
\]

If \(Q\) is terminal, then \(J(Q)={y}\), so \(y\in T_t\). If \(Q\) is nonterminal, Lemma 2 gives the controller a move \(Q\to Q'\) with

\[
J(Q')\subseteq T_t.
\]

The same argument applies when the opponent is initially first to move.

Finally, the game necessarily terminates: every move increases the number of pieces by one, while the number of pieces is at most \(n\). \(\square\)

This theorem is stronger than what is needed: either player can be designated as controller, independently of whose turn it is.

\medskip\hrule\medskip

\subsection{7. Control intervals for two initial pieces}

Suppose the current position consists of two pieces whose total length is \(N\), and write

\[
N=3q+r,\qquad r\in{0,1,2}.
\]

For residue bookkeeping, distinguish:

\begin{itemize}
\item type \(0\): length \(3i\);
\item type \(1\): length \(3i+1\);
\item type \(2_0\): the exceptional length \(2\);
\item type \(2_1\): length \(3i+2\ge5\).
\end{itemize}

Direct substitution in the definitions of \(\beta\) and \(\eta\) gives:

\begin{center}
\small
\begin{tabular}{lll}
\hline
\(N\bmod3\) & Types of the two pieces & \(J(P)\) \\
\hline
\(0\) & \(0+0\) & (\{q\}) \\
\(0\) & \(1+2_0\) & (\{q\}) \\
\(0\) & \(1+2_1\) & (\{q-1,q\}) \\
\(1\) & \(0+1\) & (\{q\}) \\
\(1\) & \(2_0+2_0\) & (\{q+1\}) \\
\(1\) & \(2_0+2_1\) & (\{q,q+1\}) \\
\(1\) & \(2_1+2_1\) & (\{q\}) \\
\(2\) & \(1+1\) & (\{q\}) \\
\(2\) & \(0+2_0\) & (\{q+1\}) \\
\(2\) & \(0+2_1\) & (\{q,q+1\}) \\
\hline
\end{tabular}
\end{center}

For example, in the last row the pieces are \(3i\) and \(3j+2\), with \(j\ge1\). Then \(q=i+j\),

\[
b(P)=i+j=q,\qquad h(P)=1,
\]

so \(J(P)={q,q+1}\). The other rows follow identically.

We therefore have the two key inclusions:

\[
\boxed{
N\equiv1,2\pmod3
\quad\Longrightarrow\quad
J(P)\subseteq{q,q+1},
}
\]

and

\[
\boxed{
N\equiv0\pmod3
\quad\Longrightarrow\quad
J(P)\subseteq{q-1,q}.
}
\]

These statements hold for every two-piece position of total length \(N\), whether or not the two lengths are distinct.

\medskip\hrule\medskip

\subsection{8. Determination of the winner}

After \(A\)'s opening move, there are two pieces and \(B\) is next.

\subsubsection{Case 1: \(n=3q+1\)}

Every opening produces a two-piece position \(P\) satisfying

\[
J(P)\subseteq{q,q+1}.
\]

Designate \(B\) as controller. By Theorem 3, \(B\) can force

\[
y\in{q,q+1}.
\]

We check both possibilities.

\paragraph{If \(y=q\)}

Then

\[
x=n-2y=3q+1-2q=q+1>q=y.
\]

Also,

\[
R=n-y-2=3q+1-q-2=2q-1
\]

is odd, so \(B\) is the last mover. Since \(x>y\), the last mover wins. Thus \(B\) wins.

\paragraph{If \(y=q+1\)}

Then

\[
x=3q+1-2(q+1)=q-1<q+1=y.
\]

Moreover,

\[
R=3q+1-(q+1)-2=2q-2
\]

is even, so \(A\) is the last mover. Since \(x<y\), the last mover loses. Thus \(B\) again wins.

Therefore

\[
\boxed{n\equiv1\pmod3\quad\Longrightarrow\quad B\text{ wins}.}
\]

\medskip\hrule\medskip

\subsubsection{Case 2: \(n=3q+2\)}

Here \(n\ge5\), since \(n=2\) lies outside the legal range.

Every opening gives a position satisfying

\[
J(P)\subseteq{q,q+1}.
\]

This time designate \(A\) as controller. Although \(B\) is next to move, Theorem 3 applies regardless of who moves next. Hence \(A\) can force

\[
y\in{q,q+1}.
\]

\paragraph{If \(y=q\)}

Then

\[
x=3q+2-2q=q+2>q=y,
\]

and

\[
R=3q+2-q-2=2q
\]

is even. Therefore \(A\) is the last mover and wins.

\paragraph{If \(y=q+1\)}

Then

\[
x=3q+2-2(q+1)=q<q+1=y,
\]

while

\[
R=3q+2-(q+1)-2=2q-1
\]

is odd. Thus \(B\) is the last mover, but because \(x<y\), the last mover loses. Hence \(A\) wins.

Therefore

\[
\boxed{n\equiv2\pmod3\quad\Longrightarrow\quad A\text{ wins}.}
\]

In fact, after any legal opening, \(A\) can use the interval-control strategy.

\medskip\hrule\medskip

\subsubsection{Case 3: \(n=3q\)}

We prove separately that \(B\) cannot lose and that \(A\) cannot lose.

\paragraph{\(B\) has a non-losing strategy}

Whatever opening \(A\) chooses, the resulting two-piece position satisfies

\[
J(P)\subseteq{q-1,q}.
\]

Designate \(B\) as controller. Then \(B\) can force

\[
y\in{q-1,q}.
\]

\begin{itemize}
\item If \(y=q\), then \(x=q\), so the game is drawn.
\item If \(y=q-1\), then
\end{itemize}
\[
  x=3q-2(q-1)=q+2>q-1=y,
\]
  and
\[
  R=3q-(q-1)-2=2q-1
\]
  is odd. Hence \(B\) is the last mover and wins.

Thus \(B\) can force either a win or a draw.

\paragraph{\(A\) has a non-losing strategy}

Player \(A\) chooses the following opening:

\[
\begin{cases}
1+2,&n=3,\\
2+4,&n=6,\\
3+(n-3),&n\ge9.
\end{cases}
\]

All these splits are legal and have distinct positive parts. Their control interval is exactly

\[
J(P)={q}.
\]

Indeed:

\begin{itemize}
\item for \(n=3\), \(b(1)+b(2)=1=q\);
\item for \(n=6\), \(b(2)+b(4)=1+1=2=q\);
\item for \(n=3q\ge9\),
\end{itemize}
\[
  b(3)+b(3q-3)=1+(q-1)=q,
\]
  and there are no hot pieces.

Now designate \(A\) as controller with target

\[
{q,q+1}.
\]

Although \(B\) moves next, Theorem 3 allows \(A\) to force

\[
y\in{q,q+1}.
\]

\begin{itemize}
\item If \(y=q\), the game is drawn.
\item If \(y=q+1\), then, whenever this value is feasible,
\end{itemize}
\[
  x=3q-2(q+1)=q-2<q+1=y,
\]
  and
\[
  R=3q-(q+1)-2=2q-3
\]
  is odd. Thus \(B\) is the last mover and loses, so \(A\) wins.

Hence \(A\) can force either a win or a draw.

Since both players possess non-losing strategies, neither can force a win against optimal opposition. Therefore

\[
\boxed{n\equiv0\pmod3\quad\Longrightarrow\quad\text{draw}.}
\]

If both players use the displayed control strategies, \(B\) forces \(y\in{q-1,q}\) while \(A\) forces \(y\in{q,q+1}\); consequently \(y=q\), so the resulting play is necessarily a draw.

\medskip\hrule\medskip

\subsection{9. Explicit optimal strategy and complexity}

The proof gives a direct algorithm.

Suppose a player has been designated as controller for a target

\[
T_t={t,t+1}.
\]

At each controller turn, compute \(b(P)\), \(h(P)\), and \(J(P)\).

\paragraph{When \(h(P)\) is odd}

Choose any hot piece \(m\).

\begin{itemize}
\item Splitting \(m=1+(m-1)\) makes the new interval the lower endpoint of \(J(P)\).
\item Splitting \(m=2+(m-2)\) makes it the upper endpoint.
\end{itemize}

Choose an endpoint lying in \(T_t\).

\paragraph{When \(h(P)>0\) is even}

Then \(J(P)={s}\), where \(s=t\) or \(s=t+1\). Choose a hot piece \(m\).

\begin{itemize}
\item If \(s=t\), use \(m=2+(m-2)\).
\item If \(s=t+1\), use \(m=1+(m-1)\).
\end{itemize}

The resulting interval is exactly \(T_t\).

\paragraph{When \(h(P)=0\)}

Choose any splittable piece \(m\).

\begin{itemize}
\item If \(m=3\), play \(1+2\).
\item If \(m\equiv0\pmod3\) and \(m\ge6\), play (2+\(m-2\)).
\item If \(m\equiv1\pmod3\), play (1+\(m-1\)).
\end{itemize}

This preserves \(J(P)\).

The appropriate controllers and targets are:

\[
\begin{array}{c|c|c}
n\bmod3 & \text{controller after the opening} & \text{target}\\ \hline
1 & B & {q,q+1}\\
2 & A & {q,q+1}\\
0 & B & {q-1,q}
\end{array}
\]

For \(n\equiv0\pmod3\), \(A\) additionally uses the special opening above and controls (\{q,q+1\}).

If the multiset of pieces is maintained together with \(b(P)\), \(h(P)\), and buckets for hot and non-hot active pieces, every update and every strategic choice takes \(O(1)\) word operations. There are at most \(n-1\) moves, since every move increases the number of pieces and there can be at most \(n\) pieces. Thus the full strategy is implementable in

\[
O(n)
\]

word operations and \(O(n)\) storage. A naive implementation that scans all current pieces at every turn has \(O(n^2)\) worst-case running time.

\medskip\hrule\medskip

\subsection{10. Boundary checks}

\paragraph{\(n=3\)}

The only move is

\[
3\longrightarrow1+2.
\]

Thus \(x=y=1\), and the game is drawn.

\paragraph{\(n=4\)}

The only opening is

\[
4\longrightarrow1+3.
\]

Then \(B\) must split \(3\to1+2\). The terminal position has \(x=2,y=1\), and \(B\) is the last mover, so \(B\) wins.

\paragraph{\(n=5\)}

Player \(A\) may play

\[
5\longrightarrow2+3.
\]

Player \(B\) must split \(3\to1+2\). The terminal pieces are (1,2,2), so \(x=1<2=y\). Since \(B\) is the last mover, \(B\) loses and \(A\) wins.

\paragraph{\(n=6\)}

Player \(A\) uses

\[
6\longrightarrow2+4.
\]

Then the continuation is forced:

\[
4\longrightarrow1+3,
\qquad
3\longrightarrow1+2.
\]

The terminal position has \(x=y=2\), hence a draw.

These are exactly the first four cases predicted by the theorem.

\medskip\hrule\medskip

\subsection{11. Necessity of the hypotheses}

\begin{itemize}
\item The restriction \(n\ge3\) is necessary merely to make the stated payoff rule well-defined: for \(n=1,2\), no one ever plays, so there is no last player.
\item The requirement that the two new lengths be \textbf{distinct} is essential. If equal splits were permitted, then for \(n=3\) one would have
\end{itemize}
\[
  3\to1+2\to1+1+1,
\]
  so \(B\), rather than producing a draw, would be the last mover with \(x>y\) and would win.
\begin{itemize}
\item Integrality is used both for termination at pieces of lengths (1,2) and for the residue-\(3\) control invariant.
\end{itemize}

\medskip\hrule\medskip

\subsection{12. Exact computational verification}

The proof above is independent of computation. The following exact dynamic program was used only as a discovery and consistency check. It verifies the claimed pattern for \(3\le n\le40\).

\begin{verbatim}
from functools import cache


def legal_moves(position: tuple[int, ...]) -> set[tuple[int, ...]]:
    """All legal successor positions."""
    successors: set[tuple[int, ...]] = set()

    for i, m in enumerate(position):
        # a < m-a, so equal splits are automatically excluded.
        for a in range(1, (m + 1) // 2):
            b = m - a
            if a < b:
                nxt = list(position)
                nxt.pop(i)
                nxt.extend((a, b))
                successors.add(tuple(sorted(nxt)))

    return successors


@cache
def value(position: tuple[int, ...]) -> int:
    """
    Minimax value for the player whose turn it is:
      +1 = that player can force a win,
       0 = that player can force a draw,
      -1 = that player must lose under optimal play.
    """
    successors = legal_moves(position)

    if not successors:
        x = position.count(1)
        y = position.count(2)

        # At a genuine terminal position, the player to move is
        # the opponent of the last mover. That player wins iff x < y.
        return (y > x) - (y < x)

    # Roles are reversed after every move.
    return max(-value(nxt) for nxt in successors)


expected = {
    0: 0,    # draw
    1: -1,   # A loses, hence B wins
    2: 1,    # A wins
}

for n in range(3, 41):
    assert value((n,)) == expected[n % 3]

print("Verified for every n from 3 through 40.")
\end{verbatim}

For a fixed \(n\), the cached verifier visits at most the \(p(n)\) integer partitions of \(n\), where \(p\) is the partition function. There are \(O(n)\) potential splits per state, so the conceptual transition count is (O(n,p\(n\))). This is only a finite verification method; the invariant proof establishes the result for all \(n\ge3\).

\medskip\hrule\medskip

\subsubsection{Final conclusion}

The answer to Q728 is \textbf{yes}, with the necessary clarification that the game is defined for \(n\ge3\):

\[
\boxed{
\begin{aligned}
n\equiv2\pmod3&\implies A\text{ wins},\\
n\equiv1\pmod3&\implies B\text{ wins},\\
n\equiv0\pmod3&\implies\text{draw}.
\end{aligned}}
\]
